% !TEX program = xelatex
% GGSO Technical Report: English, two-column, A4.
% Compile with: xelatex -interaction=nonstopmode -halt-on-error GGSO_Technical_Report_EN.tex
%
% FONT: Uses actual Times New Roman when installed (typical Windows/Overleaf
% private font configuration). Tinos is used only as a metrically compatible
% fallback in environments without Times New Roman; no font files are bundled.
%
% FIGURES: Each figure currently contains ONLY a named placeholder.
% Replace \figureplaceholder{...}{...}{...} with a figure environment and
% \includegraphics[width=\columnwidth]{figures/your_file.pdf} when ready.
\documentclass[a4paper,10pt,twocolumn]{article}

\usepackage[a4paper,top=1.65cm,bottom=1.60cm,left=1.63cm,right=1.63cm]{geometry}
\usepackage{fontspec}
\IfFontExistsTF{Times New Roman}{\setmainfont{Times New Roman}}{\setmainfont{Tinos}}
\setmonofont{DejaVu Sans Mono}[Scale=0.78]
\usepackage{amsmath,amssymb,mathtools}
\usepackage{newtxmath}
\usepackage{graphicx,booktabs,tabularx,array,multirow}
\usepackage{microtype}
\usepackage{enumitem}
\usepackage{xcolor}
\usepackage{caption}
\usepackage{float}
\usepackage{needspace}
\usepackage{titlesec}
\usepackage{fancyhdr}
\usepackage[unicode,hidelinks]{hyperref}

\definecolor{darkink}{HTML}{171717}
\definecolor{softline}{HTML}{A4A4A4}
\definecolor{linkgray}{HTML}{2A4C73}
\color{darkink}
\hypersetup{colorlinks=true,linkcolor=linkgray,urlcolor=linkgray,citecolor=linkgray}
\setlength{\columnsep}{0.58cm}
\setlength{\parindent}{0.95em}
\setlength{\parskip}{0pt}
\linespread{1.00}
\setlist[itemize]{leftmargin=*,nosep,topsep=2pt}
\setlist[enumerate]{leftmargin=*,nosep,topsep=2pt}
\titleformat{\section}{\normalfont\bfseries\fontsize{10.5}{12}\selectfont}{\thesection.}{0.48em}{}
\titleformat{\subsection}{\normalfont\bfseries\fontsize{9.6}{11}\selectfont}{\thesubsection}{0.5em}{}
\titlespacing*{\section}{0pt}{9pt}{4pt}
\titlespacing*{\subsection}{0pt}{6pt}{2pt}
\captionsetup{font=small,labelfont=bf,skip=3pt}
\renewcommand{\refname}{References}
\setlength{\textfloatsep}{8pt plus 2pt minus 2pt}
\setlength{\floatsep}{6pt plus 2pt minus 2pt}
\setlength{\intextsep}{6pt plus 2pt minus 2pt}
\pagestyle{fancy}
\fancyhf{}
\fancyfoot[C]{\footnotesize\thepage}
\renewcommand{\headrulewidth}{0pt}
\setcounter{secnumdepth}{2}
\newcommand{\sigmoid}{\operatorname{sigmoid}}
\newcommand{\logit}{\operatorname{logit}}
\newcommand{\Quant}{\operatorname{Quantile}}
\newcommand{\LN}{\operatorname{LayerNorm}}
\newcommand{\GRU}{\operatorname{GRU}}
\newcommand{\Huber}{\operatorname{Huber}}
\newcommand{\INVALID}{\textsc{invalid}}
\newcommand{\HOLD}{\textsc{hold}}
\newcommand{\Vset}{\mathcal V}
\newcommand{\Aset}{\mathcal A}
\newcommand{\Edge}{\mathcal E}
\newcommand{\given}{\,\vert\,}

% Named placeholder: no graphic, fake screenshot, or TikZ diagram is generated.
% Later replace this macro call with \includegraphics inside a normal figure.
\newcommand{\figureplaceholder}[4][2.15cm]{%
  \begin{figure}[H]
  \centering
  \fbox{\begin{minipage}[c][#1][c]{0.94\columnwidth}
    \centering\footnotesize\textbf{Figure placeholder}\\[2pt]
    \texttt{\detokenize{#2}}
  \end{minipage}}
  \caption{#3}
  \label{#4}
  \end{figure}%
}



\begin{document}
\twocolumn[{%
\begin{center}
  {\Large\bfseries GGSO: Graph-Guided Surge Optimization\\[1.5pt]
   for Network-Wide Ride-Hailing Pricing\par}

  \vspace{5pt}
  \begin{minipage}{0.94\textwidth}
  \noindent\textbf{Abstract:}
This project develops GGSO, an automated pricing system that generates
network-wide surge adjustments for ride-hailing operations. Every five
minutes, the DuoCast model forecasts booking demand and available drivers across
H3 cells. It uses spatial adjacency and observed driver-mobility networks to
capture relationships between areas. Cell-level response models then estimate
how candidate prices may affect bookings and fulfillment, while a mixed-integer
linear program selects a coordinated network pricing policy under business
and operational constraints. Validated policies are applied automatically,
while a monitoring dashboard provides execution status and a complete audit
trail.
    \par\vspace{2pt}
    \noindent\textbf{Keywords:} Ride-hailing, surge pricing, spatiotemporal GNN,
    demand elasticity, conservative optimization, MILP, automated pricing.
  \end{minipage}
\end{center}
\vspace{6pt}
}]

% ======================== MAIN CONTENT ========================
\section{Introduction and Business Problem}
Ride-hailing demand and driver availability change quickly across locations and
five-minute intervals. One H3 cell may face a shortage while a nearby cell has
excess supply. These cells are not independent because drivers move across cell
boundaries and market conditions can extend beyond a single area
\cite{geng2019,besbes2021}.

The current surge system assigns a multiplier to each H3 cell using only local
information. This approach captures the market conditions within each cell, but
it does not account for the spatial relationships among cells in a ride-hailing
network. Similar local demand-supply conditions may therefore have different
implications depending on nearby supply and observed driver movements.

GGSO builds on the local approach by incorporating information from neighboring
cells and the observed driver-mobility network. H3 adjacency represents spatial
neighborhood relationships, while directed mobility links describe how drivers
move between cells. This network context supports coordinated surge decisions
that reflect market conditions beyond individual cells.

GGSO evaluates surge multipliers jointly across the network rather than
treating each cell in isolation. Its business objective is to improve
network-level gross merchandise value (GMV) while maintaining acceptable
booking rates and valid pricing. Every five minutes, the system automatically
applies a complete network pricing policy that passes all validation checks.
If validation fails, the current effective prices remain unchanged.

\section{Proposed GGSO Framework}
GGSO consists of three core components. \textbf{DuoCast} forecasts booking
demand and available drivers for every H3 cell using local, spatial, and
driver-mobility information. \textbf{Response models} estimate how candidate
surge multipliers may change bookings and fulfillment in each cell.
\textbf{Network optimization} uses a mixed-integer linear program (MILP) to
select one candidate per cell while enforcing network and pricing constraints.

At each five-minute decision cutoff, GGSO freezes a complete network snapshot
containing historical marketplace observations, spatial and mobility graphs,
current effective surge multipliers, and the operational context required by
the response models. DuoCast runs once on this snapshot to forecast demand and
available supply for every H3 cell.

The response layer combines these forecasts with the fitted response parameters
to evaluate the eligible pricing candidates. The MILP then selects one
candidate for each cell as a coordinated network policy. An independent
validator recomputes network coverage, hard constraints, and objective
components before the policy is applied automatically. If validation fails,
the current effective prices remain unchanged. Figure~\ref{fig:system} shows
this end-to-end flow.


\begin{figure}[htbp]
    \centering
    \includegraphics[width=1\linewidth]{figures/ggso_system_architechture.png}
    \caption{End-to-end GGSO architecture: point-in-time data, forecast,
    candidate simulation, MILP, validation, automated execution, and monitoring
    dashboard.}
    \label{fig:system}
\end{figure}

\section{DuoCast Model}
\subsection{Spatiotemporal Representation}
DuoCast is a temporal, spatial, and directed-mobility GNN. For each H3 cell
$i$, a shared gated recurrent unit (GRU) encodes historical marketplace
observations. Two graph-attention branches aggregate context from spatially
adjacent cells and directed, observed physical driver movements, respectively.
A gated residual fusion combines the temporal, spatial, and mobility
representations before two separate demand and supply prediction heads
\cite{cho2014,brody2022,geng2019}.
The forecasting function is written as
\begin{equation}
 (\widehat D_i,\widehat S_i)=
 f_{\theta}(X_{t-T+1:t},G_{\mathrm{sp}},G_{\mathrm{od},t}),
 \label{eq:forecast}
\end{equation}
where $X$ is the history available at cutoff $t$; $T$ is the number of
observed snapshots; $G_{\mathrm{sp}}$ is the H3 adjacency graph;
$G_{\mathrm{od},t}$ is the directed historical driver-mobility graph;
and $\theta$ is the set of trainable parameters. The model produces two outputs
for each H3 cell: $\widehat D_i$, the forecast booking demand in $(t,t+5]$, and
$\widehat S_i$, the forecast available driver supply at $t+5$. Driver supply
$S$ is distinct from service capacity $K$.

Spatial and mobility edges encode different relationships. H3 adjacency
represents the neighboring-cell structure, whereas the mobility graph summarizes
directed historical driver movements between cells; it is not the same
as passenger trip demand and does not predict future repositioning caused
by a new surge price. A model may improve forecasts under the historical
policy and still become unreliable outside the market conditions represented
by its data.

\subsection{Multi-Task Training Objective}
Demand and supply are trained jointly with equal contributions to the loss.
Both nonnegative targets are transformed using $\log(1+x)$ to reduce scale
differences between sparse and high-volume cells. The training objective is
\begin{equation}
\begin{aligned}
 \mathcal L_{\mathrm{fc}}=\frac{1}{2N}\sum_{i=1}^{N}\bigl[&
 \rho_{\delta}(\widehat\ell_i^D-\log(1+D_i))\\
 &+\rho_{\delta}(\widehat\ell_i^S-\log(1+S_i))\bigr],
\end{aligned}
\label{eq:forecastloss}
\end{equation}
where $N$ is the number of H3 cells; $D_i$ is the observed demand over the
next five-minute interval; $S_i$ is the observed supply at the target time;
$\widehat\ell_i^D$ and $\widehat\ell_i^S$ are the corresponding log-space
predictions; and $\rho_{\delta}$ is the Huber penalty with threshold $\delta$.
The factor $1/(2N)$ averages the loss across cells and the two targets, so
neither task receives an additional weight. Predictions are decoded as
$\widehat y=\max\{\exp(\widehat\ell)-1,0\}$. The Huber penalty is defined in
Appendix~\ref{app:gnn}.

\section{Marketplace Response and Optimization}
\subsection{Cell-Level Response Models}
GGSO fits two generalized linear models (GLMs) offline and independently for
each eligible H3 cell using historical data. The demand GLM uses applied
surge, booking counts, and demand context known before the pricing decision.
The fulfillment GLM uses matched booking and completion cohorts, effective
service capacity, and its own pre-decision context. At decision time, GGSO
loads the fitted model parameters rather than refitting either model for each
candidate.

\paragraph{Demand response GLM.}
Booking demand is modeled with a Negative Binomial GLM:
\begin{equation}
\begin{aligned}
 D_{i,t}&\sim\mathrm{NegBin}(\mu_{i,t},\kappa_i),\\
 \log\mu_{i,t}&=\alpha_i+w_i^{\top}c^D_{i,t}
                   -\beta_i\log s_{i,t},
\end{aligned}
\label{eq:nb}
\end{equation}
Here $D_{i,t}$ is the observed booking count, $\mu_{i,t}$ is its expected
value, and $\kappa_i$ is the dispersion parameter. The vector $c^D_{i,t}$
contains demand context known before the pricing decision, while $s_{i,t}>0$
is the applied surge multiplier. The parameters $\alpha_i$ and $w_i$ capture the
cell baseline and context effects. The coefficient $0\le\beta_i\le4$
controls the estimated demand response to surge: a larger $\beta_i$ implies
a stronger decrease in expected bookings as surge increases. During candidate
simulation, $\beta_i$ adjusts the DuoCast demand forecast; the fitted mean
$\mu_{i,t}$ does not replace that forecast.

\paragraph{Fulfillment response GLM.}
Among the requests in the same matured booking cohort, completed rides follow
a Binomial GLM:
\begin{equation}
\begin{aligned}
 C_{i,t}&\sim\mathrm{Binomial}(D_{i,t},\pi_{i,t}),\\
 \logit(\pi_{i,t})&=b_i+v_i^{\top}c^F_{i,t}
  +\gamma_i\log\frac{K_{i,t}+\varepsilon}{D_{i,t}+\varepsilon}.
\end{aligned}
\label{eq:binomial}
\end{equation}
Here $C_{i,t}$ counts completed rides from the $D_{i,t}$ booking cohort, and
$\pi_{i,t}$ is the probability that a booking is fulfilled before the explicit
capacity cap is applied. The estimate depends on fulfillment context
$c^F_{i,t}$ and on the ratio between service capacity $K_{i,t}$ and booking
demand. The parameters $b_i$ and $v_i$ capture the cell baseline and context effects,
$\varepsilon>0$ stabilizes the ratio, and $0\le\gamma_i\le1$ controls its
influence on fulfillment. The model output is the fulfillment probability
$\pi_{i,t}$, which becomes the fulfillment response used in candidate
simulation.

The models remain separate because booking behavior and operational
fulfillment are different mechanisms. Surge enters the demand model directly,
whereas the fulfillment model responds to capacity relative to demand. Their
fitted coefficients are combined during candidate simulation, with the final
completion estimate capped by available capacity. Both models require separate
support and calibration checks. Because they are estimated from observational
data, their coefficients are not treated as established causal price effects.
Appendix~\ref{app:response} gives the likelihood objectives and fitting checks.

\subsection{Candidate Simulation and Conservative Scoring}
For each cell, GGSO applies an adjustment
$a\in\mathcal A=\{-0.3,-0.2,\ldots,0.2,0.3\}$ to the current surge $s_i^0$,
giving $s_i(a)=s_i^0+a$. Candidates outside the configured pricing bounds or
historical support are rejected. For each eligible candidate, GGSO estimates
demand, completed rides, and modeled GMV as
\begin{equation}
 \widehat D_i(a)=\widehat D_i
 \left(\frac{s_i(a)}{s_i^0}\right)^{-\beta_i},
 \label{eq:response}
\end{equation}
\begin{equation}
\begin{aligned}
 \widehat C_i(a)&=\min\{\widehat D_i(a)FR_i(a),K_i\},\\
 U_i(a)&=\widehat C_i(a)\,p_i\,s_i(a),
\end{aligned}
\label{eq:utility}
\end{equation}
where $FR_i(a)$ is the fitted fulfillment response and $p_i$ is the
reference fare before surge. Thus, $U_i(a)$ is the modeled GMV for candidate
$a$, not its improvement over the current policy.

The GLM coefficients $\beta_i$ and $\gamma_i$ are uncertain because they are
estimated from historical data. The response artifact therefore stores $B$
paired parameter scenarios
$\vartheta_i^{(b)}=(\beta_i^{(b)},\gamma_i^{(b)})$. Under each scenario, GGSO
recomputes both the candidate and the reference action using the same parameter
pair:
\begin{equation}
\begin{aligned}
 \Delta U_i^{(b)}(a)&=U_i(a;\vartheta_i^{(b)})
                     -U_i(0;\vartheta_i^{(b)}),\\
 \widetilde U_i(a)&=\Quant_{0.05}
       \left(\left\{\Delta U_i^{(b)}(a)\right\}_{b=1}^{B}\right).
\end{aligned}
 \label{eq:quantile}
\end{equation}
The fifth percentile represents lower-performing parameter scenarios and is
used as a conservative score to reduce the risk of overconfident decisions.

All candidates share one DuoCast forecast under the current policy. The
simulator changes the response calculations for each action; it does not rerun
DuoCast or predict driver redistribution. Its validated candidate table is the
input to the network optimizer.

\subsection{Joint Network MILP}
Let $z_{i,a}\in\{0,1\}$ indicate whether cell $i$ selects adjustment $a$,
with exactly one permitted candidate selected for each cell, and let $s_i$
denote the resulting surge multiplier. The network optimizer solves
\begin{equation}
\begin{aligned}
 \max_{z}\quad
 &\sum_{i,a}z_{i,a}\widetilde U_i(a)
 -\lambda_t\sum_{i,a}z_{i,a}|a|\\
 &-\lambda_g\sum_{(i,j)\in\Edge}w_{ij}|s_i-s_j|,
\end{aligned}
\label{eq:milp}
\end{equation}
The first term rewards conservative incremental GMV. The second discourages
large adjustments from the current surge, while the third discourages large
surge differences between neighboring cells. The coefficients $\lambda_t$
and $\lambda_g$ control the two penalties, and $w_{ij}$ is the weight of edge
$(i,j)$. The solver linearizes $|s_i-s_j|$ internally as described in
Appendix~\ref{app:milp}. Two principal hard constraints are
\begin{equation}
\begin{aligned}
 BR_{\min}Q&\le\sum_{i,a}z_{i,a}\widehat D_i(a)
                       \le BR_{\max}Q,\\
 \sum_{i,a}z_{i,a}\Delta U_i(a)&\ge0,
 \qquad Q=\sum_i Q_i>0,
\end{aligned}
\label{eq:constraints}
\end{equation}
where $Q_i$ is the compatible number of quote opportunities, and
$\Delta U_i(a)=U_i(a)-U_i(0)$ is nominal modeled incremental GMV.
The no-loss constraint uses nominal modeled incremental GMV, while the
objective uses the conservative lower-tail score. It prevents a negative
nominal network total but does not guarantee realized GMV. HiGHS (High
Performance Optimization Software) solves the MILP, and an independent
validator checks network coverage and all hard constraints before the
selected policy is applied. Further details are in
Appendix~\ref{app:milp}.

\begin{figure}[htbp]
    \centering
    \includegraphics[width=1\linewidth]{figures/candidate_milp_validation_flow.png}
    \caption{Candidate-to-decision flow: supported price grid, response simulation, paired uncertainty, joint MILP, and independent validity gate.}
    \label{fig:candidate-flow}
\end{figure}

\section{Validation and Evidence}
\label{sec:validation}
The evaluation assesses both the technical performance of GGSO and its
impact on business outcomes.

\subsection{Forecasting Performance}
DuoCast is evaluated on a chronologically held-out test set using raw-space
and log-space forecasting metrics. LightGBM \cite{ke2017}, GRU-only,
Spatial GATv2, and Mobility GATv2 use the same data split, forecast horizon,
node set, and target definitions.


\begin{table}[H]
\centering
\caption{Forecasting performance on the held-out test set.}
\label{tab:forecast-performance}
\footnotesize
\setlength{\tabcolsep}{3pt}
\resizebox{\columnwidth}{!}{%
\begin{tabular}{@{}lcccccc@{}}
\toprule
\textbf{Model}
& \textbf{D-MAE} $\downarrow$
& \textbf{S-MAE} $\downarrow$
& \textbf{Log D-MAE} $\downarrow$
& \textbf{Log S-MAE} $\downarrow$
& \textbf{Log D/S MAE} $\downarrow$
& \textbf{Log D/S RMSE} $\downarrow$ \\
\midrule
LightGBM baseline
& 74.2329 & 5.8163 & 0.5149 & 0.1639 & 0.4587 & 0.5794 \\

GRU-only
& 10.0524 & 4.6252 & 0.0780 & 0.1317 & 0.1489 & 0.1898 \\

Spatial GATv2
& 9.8521 & 4.6257 & 0.0768 & 0.1318 & 0.1484 & 0.1890 \\

Mobility GATv2
& 9.8627 & 4.6268 & 0.0769 & 0.1318 & 0.1484 & 0.1890 \\

DuoCast
& \textbf{9.7866}
& \textbf{4.6237}
& \textbf{0.0765}
& \textbf{0.1317}
& \textbf{0.1482}
& \textbf{0.1888} \\
\bottomrule
\end{tabular}%
}
\end{table}



Table~\ref{tab:forecast-performance} compares DuoCast with
the baseline and ablation models on the held-out test set.
DuoCast achieves the lowest demand MAE, supply MAE, log-demand MAE,
log demand-to-supply MAE, and log demand-to-supply RMSE, while matching
the GRU-only result for log-supply MAE.

Compared with the GRU-only baseline, DuoCast reduces demand MAE by 2.64\%
and supply MAE by 0.03\%. The MAE and RMSE of the log demand-to-supply ratio
are reduced by 0.47\% and 0.53\%, respectively.

Both Spatial GATv2 and Mobility GATv2 outperform the
GRU-only baseline on demand forecasting and demand-to-supply ratio
estimation, indicating that spatial adjacency and directed driver mobility
provide complementary predictive information.

\subsection{Response Model Performance}
The demand and fulfillment GLMs are evaluated independently on
chronologically held-out observations. For the demand model, prediction error
is summarized using MAE and RMSE. For the fulfillment model, fulfillment-rate
MAE measures the difference between predicted and observed fulfillment rates.
The evaluation also reports the fitted response coefficients, their uncertainty,
and the proportion of H3 cells with validated response models.

\begin{table}[H]
\centering
\caption{Negative Binomial demand GLM performance}
\label{tab:demand-glm-performance}
\footnotesize
\begin{tabularx}{\columnwidth}{@{}Xc@{}}
\toprule
\textbf{Metric} & \textbf{Value} \\
\midrule
Demand MAE $\downarrow$                 & 6.85 \\
Demand RMSE $\downarrow$                & 10.42 \\
Median $\beta$                          & 1.20 \\
Median 95\% CI Width ($\beta$)          & 0.28 \\
Eligible H3 Cells (\%) $\uparrow$      & 92.5\% \\
\bottomrule
\end{tabularx}
\end{table}

\begin{table}[H]
\centering
\caption{Binomial fulfillment GLM performance
.}
\label{tab:fulfillment-glm-performance}
\footnotesize
\begin{tabularx}{\columnwidth}{@{}Xc@{}}
\toprule
\textbf{Metric} & \textbf{Value} \\
\midrule
FR MAE $\downarrow$                     & 0.028 \\
Median $\gamma$                         & 0.65 \\
Median 95\% CI Width ($\gamma$)         & 0.18 \\
Eligible H3 Cells (\%) $\uparrow$      & 94.2\% \\
\bottomrule
\end{tabularx}
\end{table}

The median values of $\beta$ and $\gamma$ summarize the estimated demand and
fulfillment responses across eligible H3 cells. Their interval widths indicate
the uncertainty of these estimates, while eligible-cell coverage reports the
proportion of required cells with successfully fitted and validated response
models.

These metrics are evaluated against outcomes observed under historically
applied surge multipliers. They measure predictive reliability but do not
establish causal responses to previously unobserved price changes.


\subsection{Optimization and Safety Validation}
For small test networks, the MILP solution is compared with exhaustive
enumeration of all feasible candidate combinations. This comparison verifies
the objective, candidate selection, Booking Rate bounds, nominal modeled GMV
no-loss constraint, and spatial penalty.

For each full-network decision, an independent validator recomputes candidate
eligibility, network coverage, hard constraints, and objective components from
the selected candidates. Any critical validation failure marks the decision
as \INVALID, and no new pricing policy is applied.

\subsection{Real-World Business Evaluation}
A prospective switchback experiment evaluates GGSO over a sufficiently large
operational region containing multiple connected H3 cells. The entire region
alternates between the GGSO policy and the current policy across pre-specified
randomized time blocks \cite{bojinov2023}. Total realized GMV across the
experimental region is measured for each block as the primary outcome, and the
estimated treatment effect is reported with confidence intervals. Until this
experiment is completed, GGSO's business impact remains unverified.

\section{Automated Pricing and Monitoring}

Every five minutes, GGSO uses the latest marketplace data to forecast demand
and supply, evaluate surge candidates, and select a network-wide pricing policy.
A valid policy is applied automatically. If validation fails, the current
effective prices remain unchanged.

The monitoring dashboard displays the H3 map, observed and forecast demand and
supply, current, recommended, and applied surge multipliers, decision status,
validation errors, and execution history. It updates after each pricing cycle
and retains the decision and version identity of each result.

\begin{figure}[htbp]
    \centering
    \includegraphics[width=1\linewidth]{figures/operator_console.png}
    \caption{GGSO monitoring dashboard with marketplace forecasts, surge
    recommendations, execution status, and audit history.}
    \label{fig:ui}
\end{figure}

The backend checks price bounds, precision, reference versions, and idempotency
before applying a policy, then records the execution receipt. Each new cycle
uses the latest confirmed effective surge multipliers as its reference policy.

\section{Conclusion and Next Steps}
The current surge system determines a multiplier for each H3 cell from its
local market conditions. GGSO extends this approach by incorporating spatial
adjacency and directed driver mobility across cells. DuoCast forecasts demand
and available supply, the response models estimate how each cell may react to
candidate multipliers, and the MILP selects a coordinated network pricing
policy subject to business and pricing constraints.

GGSO is designed to apply a validated network-wide pricing policy automatically
every five minutes, while the monitoring dashboard provides visibility into
forecasts, pricing decisions, validation results, and execution history.
The next steps are to validate the forecasting and response models, verify the
optimization pipeline in shadow mode, and evaluate the complete system through
a real-world switchback experiment. Until that experiment is completed,
GGSO's realized business impact remains unverified. Its business value should
ultimately be judged by its causal effect on realized network GMV.

% ============================ APPENDIX ============================
\clearpage
\appendix
\begin{center}\large\bfseries APPENDIX\end{center}
\section{Notation and Data Definitions}\label{app:data}

\subsection{Decision Setting}
Let $\Vset$ be the set of H3 resolution-7 cells and
$\Edge\subseteq\Vset\times\Vset$ the pricing-adjacency edges. At decision
time $t$, GGSO uses the current surge vector $s^0$ and data available up to
that time. The forecast horizon is $H=5$ minutes, and $T$ denotes the
configurable number of historical five-minute snapshots used by DuoCast.

\begin{table}[H]
\centering
\caption{Core notation and compatible units.}
\label{tab:notation}
\footnotesize
\begin{tabularx}{\columnwidth}{@{}p{0.30\columnwidth}X@{}}
\toprule
\textbf{Symbol} & \textbf{Definition}\\
\midrule
$\mathcal V,\mathcal E$ & Set of H3 cells and pricing-adjacency edges.\\
$i,j$ & H3 cell indices.\\
$t,T,H$ & Decision cutoff, history length, forecast horizon (minutes).\\
$D_i,\widehat D_i$ & Observed and predicted confirmed valid booking requests per five-minute cohort.\\
$S_i,\widehat S_i$ & Observed and predicted online, idle, dispatch-eligible drivers (stock at timestamp).\\
$K_i$ & Effective service capacity, measured in rides per five-minute interval; not driver stock.\\
$Q_i,Q$ & Compatible quote opportunities per cell and network total, with $Q=\sum_iQ_i$.\\
$C_i,FR_i$ & Cohort-completed rides and pre-cap fulfillment probability.\\
$s_i^0,s_i(a)$ & Effective reference surge and candidate surge multiplier.\\
$a,\mathcal A_i$ & Additive surge adjustment and eligible actions for cell $i$.\\
$p_i,U_i$ & Reference fare before surge and nominal modeled GMV.\\
$\Delta U_i,\widetilde U_i$ & Nominal modeled GMV increment and conservative incremental score.\\
$G_{\mathrm{sp}},G_{\mathrm{od}}$ & Spatial adjacency and directed observed-driver-mobility graphs.\\
$z_{i,a},d_{ij}$ & MILP binary choice and auxiliary variable for $|s_i-s_j|$.\\
$\beta_i,\gamma_i$ & Fitted price sensitivity and fulfillment response coefficients.\\
\bottomrule
\end{tabularx}
\end{table}

\subsection{Business Data Definitions}
\textbf{Demand:} $D_{i,t}$ is the number of confirmed valid booking requests
created in cell $i$ during $(t-5,t]$. The forecast $\widehat D_i$ refers to
the interval $(t,t+5]$.

\textbf{Supply:} $S_{i,t}$ is the number of online, idle, and
dispatch-eligible drivers in cell $i$ at time $t$. The forecast
$\widehat S_i$ is the corresponding driver stock at $t+5$.

\textbf{Fulfillment:} $C_{i,t}$ is the number of completed rides from the
booking cohort represented by $D_{i,t}$, measured over a defined maturity
horizon, with $0\le C_{i,t}\le D_{i,t}$.

\textbf{Quote opportunities:} $Q_{i,t}$ is the number of compatible quote
opportunities for the same cell, time window, and booking cohort.

All input features, prices, graph observations, and context variables must be
available by cutoff $t$ and aligned to the same decision interval.

\section{Forecasting Network and Training Mathematics}\label{app:gnn}
\subsection{Temporal, Spatial, and Mobility Encoders}
Let $x_{i,k}\in\mathbb R^{F}$ be the observed feature vector for cell $i$
at historical step $k$. The shared temporal encoder computes
\begin{equation}
 h_{i,k}=\GRU(x_{i,k},h_{i,k-1}),\qquad
 \ell_i=\LN(h_{i,T}).
 \label{eq:gru}
\end{equation}
The spatial relation uses H3 neighborhood edges; the mobility relation
uses physical driver transitions directed from origin to destination and
available before cutoff. Each relation $r\in\{\mathrm{sp},\mathrm{od}\}$
has independent attention parameters. For observed edge attributes
$q_{ji}^{(r)}$, one relation-aware GATv2 formulation is
\begin{equation}
\begin{aligned}
 e_{ji}^{(r)}=&\; a_r^\top \operatorname{LeakyReLU}
 \bigl(W_s^{(r)}\ell_j+W_t^{(r)}\ell_i\\
 &\hspace{42pt}+W_e^{(r)}q_{ji}^{(r)}\bigr),\\
 \alpha_{ji}^{(r)}=&\;
 \frac{\exp(e_{ji}^{(r)})}{\sum_{k\in\mathcal N_r(i)}\exp(e_{ki}^{(r)})},\\
 m_i^{(r)}=&\;\sum_{j\in\mathcal N_r(i)}
 \alpha_{ji}^{(r)}W_v^{(r)}\ell_j,
\end{aligned}
\label{eq:attention}
\end{equation}
where $\mathcal N_r(i)$ is the observed incoming-neighbor set under
relation $r$, $e_{ji}^{(r)}$ is an attention logit,
$\alpha_{ji}^{(r)}$ is its normalized coefficient, and $m_i^{(r)}$ is
the aggregated relation-specific message. Separate parameters allow the two
branches to learn different meanings for spatial proximity and observed
driver movement.

\subsection{Gated Fusion and Prediction Heads}
The temporal representation and the messages from the two graph branches are
combined through gated residual fusion:
\begin{equation}
\begin{aligned}
 c_i^{\mathrm{net}}=&\;W_c[m_i^{(\mathrm{sp})}\Vert
                         m_i^{(\mathrm{od})}],\\
 \omega_i=&\;\sigmoid\bigl(W_g[\ell_i\Vert m_i^{(\mathrm{sp})}
                  \Vert m_i^{(\mathrm{od})}]\bigr),\\
 z_i=&\;\LN(\ell_i+\omega_i\odot c_i^{\mathrm{net}}),
\end{aligned}
\label{eq:fusion}
\end{equation}
where $\Vert$ denotes concatenation, $\odot$ denotes element-wise
multiplication, and $\omega_i$ is a learned gate. Two separate heads output
$\widehat\ell_i^D=f_D(z_i)$ and $\widehat\ell_i^S=f_S(z_i)$, subsequently
decoded into nonnegative demand and supply forecasts. The gate controls how
much network context is added to the temporal representation for each cell.

\subsection{Log Transformation and Huber Loss}
Demand and supply targets are transformed using
$\widetilde y=\log(1+y)$ before training. This transformation is defined
for zero-valued observations and compresses the difference between
low-volume and high-volume H3 cells.

The prediction error is measured in log space using the Huber loss:
\begin{equation}
\rho_\delta(e)=
\begin{cases}
 \frac12e^2, & |e|\le\delta,\\
 \delta\left(|e|-\frac12\delta\right), & |e|>\delta.
\end{cases}
\label{eq:huber}
\end{equation}
where $e$ is the difference between the predicted and observed log-transformed
target, and $\delta$ controls the transition between the two regions. The loss
is quadratic for small errors and linear for large errors, reducing the
influence of unusually large demand or supply observations.

\section{Marketplace Response Estimation and Capacity}\label{app:response}
\subsection{Demand Negative Binomial Likelihood}
The expected booking demand in Eq.~\eqref{eq:nb} is
$\mu_{i,t}=\exp(\alpha_i+w_i^\top c^D_{i,t}-\beta_i\log s_{i,t})$.
The constrained maximum-likelihood training objective is
\begin{equation}
 \mathcal L_{D,i}=
 -\sum_{t\in\mathcal T_i^{\mathrm{train}}}
 \log P_{\mathrm{NB}}(D_{i,t}\mid\mu_{i,t},\kappa_i),
 \quad0\le\beta_i\le4.
 \label{eq:nbnll}
\end{equation}
The parameter tuple includes intercept $\alpha_i$, ordered context
coefficients $w_i$, elasticity $\beta_i$, and dispersion $\kappa_i$.
The Negative Binomial mean--dispersion convention must remain consistent
through fitting and evaluation because alternative parameterizations assign
different meanings to $\kappa_i$. Observed surge variation and sufficient
per-cell data are prerequisites for meaningful estimation.

\subsection{Cohort-Level Binomial Likelihood}
For Eq.~\eqref{eq:binomial}, the conditional fulfillment likelihood gives the
following negative-log-likelihood objective, with the parameter-independent
binomial coefficient omitted:
\begin{equation}
\begin{aligned}
 \mathcal L_{F,i}=-\sum_{t\in\mathcal T_i^{\mathrm{train}}}
 \bigl[&C_{i,t}\log\pi_{i,t}\\
 &+(D_{i,t}-C_{i,t})\log(1-\pi_{i,t})\bigr],
\end{aligned}
\label{eq:binnll}
\end{equation}
subject to $0\le\gamma_i\le1$ and coherent cohort counting.
When $D=0,C=0$, the interval has no positive binomial trials; it cannot
be treated as evidence for $100\%$ fulfillment. The coefficient $\gamma_i$
does not define the response by itself; the baseline $b_i$, context effects
$v_i$, and capacity-to-demand ratio are also required.

\subsection{Available Service Capacity}
If a direct, validated capacity estimate is unavailable, the proposed
approximation is
\begin{equation}
 K_i=\widehat S_i\,r_i\frac{H}{60}
 \min\!\left\{1,\frac{H}{\max(ETA_i,H)}\right\},\quad H=5,
 \label{eq:capacity}
\end{equation}
where $r_i$ is the assumed average rides per driver per hour and $ETA_i$ is a
pickup-time estimate in minutes. Thus, $K_i$ is measured in rides per five
minutes, while $\widehat S_i$ is measured in drivers. The factor $H/60$
converts the hourly service rate to the forecast horizon, and the final factor
reduces capacity when pickup time exceeds that horizon.
Equation~\eqref{eq:capacity} is a heuristic capacity conversion, not
a third trained prediction head. Unit consistency, zero-capacity handling,
and calibration require separate validation.

\subsection{Interpretation and Causal Limitations}
Because the response models are fitted on observational data, the estimated
$\beta_i$ and $\gamma_i$ are interpreted as predictive associations rather
than causal effects. Response parameters are estimated independently for
every required H3 cell.

\section{Candidate Evaluation and Conservative Scoring}\label{app:candidate}
\subsection{Candidate Evaluation}
For each H3 cell, GGSO evaluates the fixed action grid
\begin{equation}
\mathcal A=\{-0.3,-0.2,-0.1,0,0.1,0.2,0.3\},
\qquad
s_i(a)=s_i^0+a.
\label{eq:grid}
\end{equation}
A candidate is eligible when its multiplier is positive, remains within
the configured surge and adjustment bounds, and is supported by historical
prices.

For each eligible action, expected demand is
\begin{equation}
\widehat D_i(a)
=
\widehat D_i
\left(\frac{s_i(a)}{s_i^0}\right)^{-\beta_i}.
\end{equation}

The baseline fulfillment logit and candidate fulfillment probability are
\begin{equation}
\begin{aligned}
\eta_i^0
&=
b_i+v_i^\top c_i^F+
\gamma_i\log
\frac{K_i+\varepsilon}{\widehat D_i+\varepsilon},\\
FR_i(a)
&=
\sigmoid\left(
\eta_i^0+
\gamma_i\log
\frac{\widehat D_i+\varepsilon}
     {\widehat D_i(a)+\varepsilon}
\right).
\end{aligned}
\end{equation}

Expected completed rides and GMV are then
\begin{equation}
\begin{aligned}
\widehat C_i(a)
&=
\min\{\widehat D_i(a)FR_i(a),K_i\},\\
U_i(a)
&=
\widehat C_i(a)p_i s_i(a).
\end{aligned}
\end{equation}

Capacity, fulfillment context, and pre-surge fare remain fixed across all
candidates of the same cell and decision interval.

\subsection{Conservative Scoring}
Let $\vartheta_i^{(b)}=(\beta_i^{(b)},\gamma_i^{(b)})$ denote a
response-parameter scenario. The incremental GMV of action $a$ under
scenario $b$ is
\begin{equation}
\Delta U_i^{(b)}(a)
=
U_i(a;\vartheta_i^{(b)})
-
U_i(0;\vartheta_i^{(b)}).
\end{equation}

The conservative score is the empirical fifth percentile of the paired
incremental values:
\begin{equation}
\widetilde U_i(a)
=
\Quant_{0.05}
\left(\{\Delta U_i^{(b)}(a)\}_b\right).
\end{equation}

The candidate and baseline are recomputed using the same parameter scenario,
so the score measures their paired difference under response uncertainty.
For the zero action, both the incremental GMV and conservative score are zero.

\section{MILP Formulation and Validation}\label{app:milp}
\subsection{Decision Variables and Objective}
For each cell $i\in\Vset$ and eligible candidate
$a\in\mathcal A_i$, let $z_{i,a}$ indicate whether the candidate is
selected. Exactly one candidate is selected for each cell:
\begin{equation}
\begin{gathered}
 z_{i,a}\in\{0,1\},\qquad
 \sum_{a\in\mathcal A_i}z_{i,a}=1,\\
 s_i=\sum_{a\in\mathcal A_i}z_{i,a}s_i(a).
\end{gathered}
\label{eq:milpvars}
\end{equation}

For each pricing edge $(i,j)$, the auxiliary variable $d_{ij}$ linearizes
the absolute difference between the selected surge multipliers:
\begin{equation}
 d_{ij}\ge0,\qquad
 d_{ij}\ge s_i-s_j,\qquad
 d_{ij}\ge s_j-s_i.
\end{equation}
The objective maximizes conservative incremental utility while penalizing
large changes from the current price and large price differences between
neighboring cells, as defined in Eq.~\eqref{eq:milp}.

\subsection{Hard Constraints}
The predicted network Booking Rate is constrained by
\begin{equation}
 BR_{\min}Q
 \le
 \sum_{i\in\Vset}\sum_{a\in\mathcal A_i}
 z_{i,a}\widehat D_i(a)
 \le
 BR_{\max}Q,
 \qquad Q=\sum_i Q_i>0.
\end{equation}

The selected policy must also satisfy the nominal modeled GMV no-loss
constraint:
\begin{equation}
 \sum_{i\in\Vset}\sum_{a\in\mathcal A_i}
 z_{i,a}\,[U_i(a)-U_i(0)]\ge0.
 \label{eq:noloss}
\end{equation}
The feasible candidate sets already enforce the action grid, surge bounds,
adjustment limits, and historical support requirements.

\subsection{Solver and Independent Validation}
HiGHS solves the MILP. A proven optimal solution or a feasible incumbent
returned at the time limit may be accepted only if an independent validator
recomputes network coverage, candidate eligibility, Booking Rate, nominal
modeled GMV, and the objective from the selected candidates. This independent
check prevents a solver status alone from being treated as evidence that the
resulting network policy is valid. If no feasible solution exists or validation
fails, GGSO produces no recommendation and the current prices remain unchanged.

\section{Model Assumptions and Limitations}\label{app:limitations}
DuoCast assumes that historical temporal, spatial, and mobility patterns
remain informative for the next five-minute interval. Forecast reliability
may decrease when marketplace conditions differ substantially from the
training data.

The response models are estimated from observational data. Their predictive
performance does not establish that the estimated price effects are causal.

Surge pricing can induce spatial supply responses \cite{besbes2021}. GGSO
candidate simulation, however, holds capacity, fare, and pre-decision context
fixed across price candidates and does not model driver redistribution.
Cell-level conservative scores do not form a joint network-level confidence
bound, and modeled no-loss does not guarantee realized no-loss.

\begin{thebibliography}{99}
\small

\bibitem{geng2019}
X. Geng, Y. Li, L. Wang, L. Zhang, Q. Yang, J. Ye, and Y. Liu,
``Spatiotemporal Multi-Graph Convolution Network for Ride-Hailing Demand
Forecasting,'' in \textit{Proceedings of the AAAI Conference on Artificial
Intelligence}, vol. 33, no. 1, pp. 3656--3663, 2019.
doi: \href{https://doi.org/10.1609/aaai.v33i01.33013656}{10.1609/aaai.v33i01.33013656}.

\bibitem{besbes2021}
O. Besbes, F. Castro, and I. Lobel, ``Surge Pricing and Its Spatial Supply
Response,'' \textit{Management Science}, vol. 67, no. 3, pp. 1350--1367,
2021. doi:
\href{https://doi.org/10.1287/mnsc.2020.3622}{10.1287/mnsc.2020.3622}.

\bibitem{cho2014}
K. Cho, B. van Merri\"enboer, C. Gulcehre, D. Bahdanau, F. Bougares,
H. Schwenk, and Y. Bengio, ``Learning Phrase Representations using RNN
Encoder--Decoder for Statistical Machine Translation,'' in
\textit{Proceedings of EMNLP}, pp. 1724--1734, 2014.
doi: \href{https://doi.org/10.3115/v1/D14-1179}{10.3115/v1/D14-1179}.

\bibitem{brody2022}
S. Brody, U. Alon, and E. Yahav, ``How Attentive Are Graph Attention
Networks?'' in \textit{Proceedings of ICLR}, 2022.
Available: \url{https://openreview.net/forum?id=F72ximsx7C1}.

\bibitem{ke2017}
G. Ke, Q. Meng, T. Finley, T. Wang, W. Chen, W. Ma, Q. Ye, and T.-Y. Liu,
``LightGBM: A Highly Efficient Gradient Boosting Decision Tree,'' in
\textit{Advances in Neural Information Processing Systems}, vol. 30, 2017.

\bibitem{bojinov2023}
I. Bojinov, D. Simchi-Levi, and J. Zhao, ``Design and Analysis of Switchback
Experiments,'' \textit{Management Science}, vol. 69, no. 7,
pp. 3759--3777, 2023. doi:
\href{https://doi.org/10.1287/mnsc.2022.4583}{10.1287/mnsc.2022.4583}.

\end{thebibliography}

\end{document}
